\begin{center}
\begin{tikzpicture}[x=.92cm,y=.92cm,font=\small]
  \fill[paper] (0,0) rectangle (10,5.8);
  \fill[dynasty!22] (.6,1.0) -- (3.7,.8) -- (4.3,4.7) -- (1.2,5.0) -- cycle;
  \node[align=center] at (2.35,2.9) {秦\\关中};
  \fill[green!20] (5.3,1.1) -- (9.3,1.0) -- (9.2,5.0) -- (5.7,5.1) -- cycle;
  \node[align=center] at (7.4,3.25) {赵\\邯郸方向};
  \draw[gray!70,line width=1.5pt] (4.7,.3) .. controls (4.2,2.5) and (4.8,4.1) .. (4.5,5.5);
  \node[rotate=90,text=gray!70] at (4.15,3.0) {太行山地};
  \fill[timeline] (5.0,2.9) circle (.13);
  \node[above,align=center] at (5.0,3.05) {长平\\上党一带};
  \draw[-{Stealth[length=2mm]},ultra thick,color=dynasty] (2.8,2.5) -- (4.85,2.9);
  \draw[-{Stealth[length=2mm]},thick,color=green!60!black] (7.3,2.2) -- (5.15,2.85);
  \node[text=dynasty,below] at (3.8,2.3) {秦军东进与补给};
  \node[text=green!60!black,above] at (6.4,2.25) {赵军西向防御};
\end{tikzpicture}

\smallskip
\small\textit{图  长平战场空间示意：突出关中、太行、上党/长平与赵国腹地之间的补给和地形关系；不表示精确行军路线或国界。}
\end{center}
